\chapter{L'excrétion urinaire et l'insuffisance rénale}

\noindent\textit{Le fonctionnement des cellules produit des déchets qui doivent être
éliminés pour ne pas empoisonner l'organisme. C'est le rôle de l'appareil urinaire, et en
particulier des reins, qui filtrent le sang en continu pour produire l'urine. Ce chapitre
étudie leur fonctionnement, puis ce qui se passe lorsqu'ils cessent de fonctionner
correctement.}
\vspace{8pt}

\connecteBanner

\section{L'appareil urinaire}

\begin{definition}[Organes de l'appareil urinaire]
L'appareil urinaire comprend deux \textbf{reins} (qui filtrent le sang et produisent
l'urine), deux \textbf{uretères} (qui conduisent l'urine des reins à la vessie), la
\textbf{vessie} (qui stocke l'urine) et l'\textbf{urètre} (qui évacue l'urine vers
l'extérieur).
\end{definition}

\begin{illus}
\begin{tikzpicture}[scale=0.8]
\filldraw[onbo!20,draw=onbo,line width=1.1pt] (0,0) rectangle (1.9,0.9);
\node[font=\small] at (0.95,0.45) {reins};
\draw[->,onbmuted,line width=1.1pt] (1.9,0.45) -- (3.5,0.45);
\filldraw[onbgreen!20,draw=onbgreen,line width=1.1pt] (3.5,0) rectangle (5.9,0.9);
\node[font=\small] at (4.7,0.45) {uretères};
\draw[->,onbmuted,line width=1.1pt] (5.9,0.45) -- (7.5,0.45);
\filldraw[onbo2!20,draw=onbo2,line width=1.1pt] (7.5,0) rectangle (9.7,0.9);
\node[font=\small] at (8.6,0.45) {vessie};
\draw[->,onbmuted,line width=1.1pt] (9.7,0.45) -- (11.3,0.45);
\filldraw[onbink!15,draw=onbink,line width=1.1pt] (11.3,0) rectangle (13.4,0.9);
\node[font=\small] at (12.35,0.45) {urètre};
\end{tikzpicture}
\fig{Figure — Trajet de l'urine : des reins à l'extérieur du corps.}
\end{illus}

\section{Le rein, filtre du sang}

\begin{propriete}[Rôle du rein]
Le sang arrive au rein chargé de \textbf{déchets} produits par les cellules (urée,
acide urique) et d'un excès d'\textbf{eau} et de \textbf{sels minéraux}. Le rein
\textbf{filtre} en continu ce sang pour en extraire ces substances, qui forment
l'\textbf{urine}, tandis que le sang \textbf{purifié} repart dans la circulation.
\end{propriete}

\begin{definition}[Néphron]
Le \textbf{néphron} est l'unité fonctionnelle du rein (chaque rein en contient environ un
million). Il assure deux étapes successives : une \textbf{filtration} du sang (qui laisse
passer l'eau, les sels et les déchets, mais retient les cellules sanguines et les grosses
molécules comme les protéines), puis une \textbf{réabsorption} qui récupère l'essentiel de
l'eau et des substances utiles (glucose, sels) vers le sang, ne laissant dans l'urine
définitive que les déchets et le surplus d'eau et de sels.
\end{definition}

\begin{propriete}[Un rendement de filtration impressionnant]
Chaque jour, les reins filtrent environ $180$ L de liquide, mais n'excrètent qu'environ
$1{,}8$ L d'urine : plus de $99\,\%$ du liquide filtré est donc \textbf{réabsorbé} et
renvoyé dans le sang.
\end{propriete}

\begin{exemple}
Volume filtré $V_f=180$ L, volume d'urine $V_u=1{,}8$ L : pourcentage réabsorbé
$=\dfrac{V_f-V_u}{V_f}\times100=\dfrac{180-1{,}8}{180}\times100\approx99\,\%$.
\end{exemple}

\section{Le rein et la régulation du milieu intérieur}

\begin{propriete}[Homéostasie]
En ajustant la quantité d'eau et de sels minéraux éliminés dans l'urine selon les besoins
de l'organisme, le rein maintient \textbf{constante} la composition du sang (quantité
d'eau, de sels, de déchets) : c'est un exemple d'\textbf{homéostasie}, essentielle au bon
fonctionnement de toutes les cellules.
\end{propriete}

\begin{remarque}
Par forte chaleur ou après une transpiration importante, les reins produisent une urine
plus \textbf{concentrée} et en plus faible quantité, pour limiter la perte d'eau de
l'organisme.
\end{remarque}

\section{Analyser une urine}

\begin{propriete}[Composition normale et anomalies]
Une urine normale contient de l'eau, de l'urée, des sels minéraux, mais \textbf{pas} de
glucose ni de protéines en quantité notable. La présence anormale de certaines substances
oriente vers un problème de santé :
\begin{itemize}
\item \textbf{glucose dans l'urine} : évoque un \textbf{diabète} (le rein n'arrive plus à
tout réabsorber lorsque le taux de glucose sanguin est trop élevé) ;
\item \textbf{protéines dans l'urine} : évoque un problème de \textbf{filtration
rénale} ;
\item \textbf{sang dans l'urine} : évoque une \textbf{infection} ou une lésion des voies
urinaires.
\end{itemize}
\end{propriete}

\section{L'insuffisance rénale}

\begin{definition}[Insuffisance rénale]
L'\textbf{insuffisance rénale} est une diminution, progressive ou brutale, de la
capacité des reins à filtrer correctement le sang. Les déchets et l'excès d'eau et de
sels ne sont plus éliminés efficacement et s'accumulent dans l'organisme.
\end{definition}

\begin{propriete}[Causes principales]
L'hypertension artérielle non contrôlée et le diabète mal équilibré (voir chapitre
11) sont les deux principales causes d'insuffisance rénale : ils abîment
progressivement les vaisseaux sanguins des reins. Des infections rénales répétées
peuvent aussi endommager les reins.
\end{propriete}

\begin{propriete}[Symptômes]
Fatigue, gonflements (œdèmes) du visage ou des jambes, diminution du volume
d'urine, urines anormales, et parfois aggravation d'une hypertension déjà présente.
L'insuffisance rénale peut évoluer \textbf{silencieusement} pendant longtemps avant
d'être détectée.
\end{propriete}

\begin{propriete}[Prévention et prise en charge]
\begin{itemize}
\item contrôler régulièrement sa \textbf{tension artérielle} et son \textbf{taux de
glucose} sanguin ;
\item boire suffisamment d'eau et limiter une consommation excessive de sel ;
\item éviter l'\textbf{automédication} abusive (certains médicaments sont toxiques
pour les reins à forte dose) ;
\item réaliser un dépistage régulier (analyse d'urine, prise de tension) ;
\item en cas d'insuffisance rénale sévère, la \textbf{dialyse} (filtration
artificielle du sang par une machine) ou la \textbf{greffe de rein} permettent de
suppléer ou remplacer la fonction des reins défaillants.
\end{itemize}
\end{propriete}

\begin{exemple}
Une personne hypertendue qui ne suit pas de traitement pendant plusieurs années
expose progressivement ses reins à une insuffisance rénale, souvent sans aucun
symptôme avant un stade avancé.
\end{exemple}

\begin{remarque}
Piège fréquent : l'insuffisance rénale n'est pas seulement une maladie du rein
lui-même ; elle est le plus souvent la \textbf{conséquence} d'une autre maladie mal
contrôlée (hypertension, diabète), ce qui explique l'importance de prévenir et de
suivre ces maladies.
\end{remarque}

% ═══════════════════════════════════════════════════════════════
\section{L'essentiel à retenir}

\begin{synthese}
\textbf{Appareil urinaire.} Reins $\to$ uretères $\to$ vessie $\to$ urètre.\\[4pt]
\textbf{Néphron.} Filtration du sang, puis réabsorption de l'eau et des substances
utiles.\\[4pt]
\textbf{Rendement.} Environ $99\,\%$ du liquide filtré est réabsorbé.\\[4pt]
\textbf{Homéostasie.} Le rein maintient constante la composition du sang.\\[4pt]
\textbf{Anomalies de l'urine.} Glucose $\to$ diabète ; protéines $\to$ problème de
filtration ; sang $\to$ infection.\\[4pt]
\textbf{Insuffisance rénale.} Diminution de la capacité de filtration des reins,
causée surtout par l'hypertension et le diabète mal contrôlés ; traitée par dialyse
ou greffe si sévère.\\[4pt]
\textbf{Piège fréquent.} Le rein ne fabrique pas les déchets : il les élimine. Les
déchets (urée) proviennent de l'activité des cellules. L'insuffisance rénale est
souvent la conséquence d'une autre maladie mal contrôlée.
\end{synthese}

% ═══════════════════════════════════════════════════════════════
\section{Méthodes \& savoir-faire}

\methode{Retracer le trajet de l'urine}
Suivre l'ordre : reins, uretères, vessie, urètre.
\begin{exemple}
L'urine ne stagne pas dans le rein : elle est continuellement évacuée vers la vessie via
les uretères.
\end{exemple}

\methode{Calculer un pourcentage de réabsorption rénale}
Appliquer $\%\text{réabsorbé}=\dfrac{V_f-V_u}{V_f}\times100$, avec $V_f$ le volume
filtré et $V_u$ le volume d'urine excrétée.
\begin{exemple}
$V_f=150$ L, $V_u=1{,}5$ L : $\%=\dfrac{150-1{,}5}{150}\times100=99\,\%$.
\end{exemple}

\methode{Interpréter une analyse d'urine}
Repérer la substance anormalement présente et l'associer à sa cause probable : glucose
$\to$ diabète ; protéines $\to$ filtration ; sang $\to$ infection.
\begin{exemple}
Une bandelette urinaire positive au glucose oriente vers un dépistage du diabète.
\end{exemple}

\methode{Relier un facteur de risque à l'insuffisance rénale}
Identifier si la situation évoque une maladie qui abîme les reins sur le long terme
(hypertension, diabète mal contrôlé) ou un geste de prévention (hydratation,
dépistage, éviter l'automédication).
\begin{exemple}
Une hypertension artérielle non traitée pendant des années est un facteur de risque
d'insuffisance rénale.
\end{exemple}

% ═══════════════════════════════════════════════════════════════
\section{Exercices d'application}
\begin{multicols}{2}

\rubrique{Appareil urinaire}
\exo{}\diff{1} Citer, dans l'ordre, les quatre organes traversés par l'urine.

\exo{}\diff{1} Quel est le rôle de la vessie ?

\rubrique{Le rein, filtre du sang}
\exo{}\diff{2} Que fait le rein au sang qui le traverse ?

\exo{}\diff{2} Qu'est-ce qu'un néphron ?

\exo{}\diff{2} Quelles sont les deux étapes réalisées par le néphron ?

\exo{}\diff{2} Un rein filtre $V_f=150$ L de liquide par jour et excrète $V_u=1{,}5$ L
d'urine. Calculer le pourcentage réabsorbé.

\rubrique{Régulation}
\exo{}\diff{2} Qu'est-ce que l'homéostasie ?

\exo{}\diff{2} Pourquoi l'urine devient-elle plus concentrée par forte chaleur ?

\rubrique{Analyser une urine}
\exo{}\diff{2} Que peut indiquer la présence de glucose dans l'urine ?

\exo{}\diff{2} Que peut indiquer la présence de sang dans l'urine ?

\rubrique{Insuffisance rénale}
\exo{}\diff{1} Définir l'insuffisance rénale.

\exo{}\diff{2} Citer les deux principales causes d'insuffisance rénale.

\exo{}\diff{2} Citer deux mesures de prévention de l'insuffisance rénale.

% ═══════════════════════════════════════════════════════════════
\end{multicols}

\section{Exercices d'entraînement}
\begin{multicols}{2}

\rubrique{Appareil urinaire}
\exo{}\diff{2} Pourquoi l'appareil urinaire comporte-t-il deux reins mais un seul
urètre ?

\rubrique{Le rein, filtre du sang}
\exo{}\diff{3} Expliquer pourquoi les grosses molécules comme les protéines ne se
retrouvent normalement pas dans l'urine.

\exo{}\diff{2} Un rein filtre $V_f=160$ L par jour et excrète $V_u=1{,}6$ L d'urine.
Calculer le pourcentage réabsorbé.

\exo{}\diff{3} Un rein filtre $V_f=120$ L par jour et réabsorbe $99\,\%$ de ce volume.
Calculer le volume d'urine excrété.

\rubrique{Régulation}
\exo{}\diff{3} Expliquer, en une phrase, pourquoi le rein est essentiel au maintien
d'une composition sanguine stable.

\exo{}\diff{2} Une personne boit une grande quantité d'eau. Que va probablement faire son
rein en réaction ?

\rubrique{Analyser une urine}
\exo{}\diff{3} Un patient présente du glucose dans ses urines. Expliquer, à l'aide du
fonctionnement du néphron, pourquoi cela peut se produire lorsque le taux de glucose
sanguin est trop élevé.

\exo{}\diff{2} Pourquoi une bandelette urinaire est-elle un outil de dépistage rapide et
utile ?

\exo{}\diff{3} Expliquer la différence entre un déchet normalement présent dans l'urine
(comme l'urée) et une substance anormale (comme le glucose).

\rubrique{Insuffisance rénale}
\exo{}\diff{3} Expliquer pourquoi une hypertension artérielle non traitée peut, à long
terme, provoquer une insuffisance rénale.

\exo{}\diff{2} Pourquoi l'insuffisance rénale est-elle parfois découverte
tardivement ?

\exo{}\diff{3} Expliquer pourquoi la dialyse est nécessaire en cas d'insuffisance
rénale sévère.

% ═══════════════════════════════════════════════════════════════
\end{multicols}

\section{Sujets type BEPC}

\sujetbac[Le rein et la formation de l'urine]
\begin{enumerate}[label=\arabic*)]
\item Cite, dans l'ordre, les organes traversés par l'urine, du rein à l'extérieur du
corps.
\item Quel est le rôle du néphron ?
\item Un rein filtre $V_f=180$ L de liquide par jour et excrète $V_u=1{,}8$ L d'urine.
Calcule le pourcentage réabsorbé.
\item Que deviennent, en majorité, l'eau et les substances utiles filtrées par le
rein ?
\end{enumerate}

\sujetbac[Homéostasie et régulation]
\begin{enumerate}[label=\arabic*)]
\item Définis l'homéostasie.
\item Explique pourquoi l'urine est plus concentrée par forte chaleur.
\item Pourquoi le maintien d'une composition sanguine stable est-il important pour les
cellules de l'organisme ?
\item Cite un organe, autre que le rein, qui participe au maintien du milieu intérieur
(indice : chapitre sur la respiration).
\end{enumerate}

\sujetbac[Analyse d'urine et dépistage]
\begin{enumerate}[label=\arabic*)]
\item Cite trois substances qui ne doivent normalement pas être présentes en quantité
notable dans l'urine.
\item Que peut indiquer la présence de glucose dans l'urine ?
\item Que peut indiquer la présence de protéines dans l'urine ?
\item Pourquoi l'analyse d'urine est-elle un outil de dépistage utile en médecine ?
\end{enumerate}

\sujetbac[L'insuffisance rénale]
\begin{enumerate}[label=\arabic*)]
\item Définis l'insuffisance rénale.
\item Cite les deux principales causes de l'insuffisance rénale.
\item Cite trois symptômes possibles de l'insuffisance rénale.
\item Cite deux mesures de prévention et un traitement possible en cas
d'insuffisance rénale sévère.
\end{enumerate}

% ═══════════════════════════════════════════════════════════════
\section{Corrigés}
\begin{multicols}{2}

\corrige{de l'exercice 1}
Reins, uretères, vessie, urètre.

\corrige{de l'exercice 2}
Elle stocke l'urine avant son évacuation.

\corrige{de l'exercice 3}
Il le filtre, pour en extraire les déchets, l'excès d'eau et de sels, formant l'urine.

\corrige{de l'exercice 4}
L'unité fonctionnelle du rein, qui réalise la filtration du sang puis la réabsorption.

\corrige{de l'exercice 5}
La filtration, puis la réabsorption.

\corrige{de l'exercice 6}
$\dfrac{150-1{,}5}{150}\times100=99\,\%$.

\corrige{de l'exercice 7}
Le maintien constant de la composition du milieu intérieur (sang) de l'organisme.

\corrige{de l'exercice 8}
Car le rein réabsorbe davantage d'eau pour limiter les pertes, produisant une urine plus
concentrée en moins grande quantité.

\corrige{de l'exercice 9}
Un possible diabète.

\corrige{de l'exercice 10}
Une infection ou une lésion des voies urinaires.

\corrige{de l'exercice 11}
Une diminution de la capacité des reins à filtrer correctement le sang.

\corrige{de l'exercice 12}
L'hypertension artérielle non contrôlée et le diabète mal équilibré.

\corrige{de l'exercice 13}
Par exemple : contrôler sa tension artérielle et éviter l'automédication abusive.

\corrige{de l'exercice 14}
Car un seul rein suffit à assurer une filtration correcte (redondance protectrice),
mais l'urine produite par les deux reins doit être évacuée par un unique conduit
commun vers l'extérieur.

\corrige{de l'exercice 15}
Car le filtre du néphron retient normalement les grosses molécules comme les
protéines, qui restent donc dans le sang.

\corrige{de l'exercice 16}
$\dfrac{160-1{,}6}{160}\times100=99\,\%$.

\corrige{de l'exercice 17}
$V_u=120\times(1-0{,}99)=120\times0{,}01=1{,}2$ L.

\corrige{de l'exercice 18}
Car il élimine sélectivement les déchets et ajuste les quantités d'eau et de sels
excrétées, maintenant une composition sanguine stable indispensable au bon
fonctionnement des cellules.

\corrige{de l'exercice 19}
Il va produire une urine plus abondante et plus diluée, pour éliminer l'excès d'eau
absorbée.

\corrige{de l'exercice 20}
Car une partie du glucose filtré n'est plus totalement réabsorbée par le néphron
lorsque sa quantité dans le sang est trop élevée ; le glucose en excès se retrouve
alors dans l'urine.

\corrige{de l'exercice 21}
Car elle donne un résultat rapide, simple et peu coûteux, permettant d'orienter
ensuite vers des examens plus poussés si nécessaire.

\corrige{de l'exercice 22}
L'urée est un déchet normal, produit en permanence par l'activité des cellules et
toujours présent dans l'urine ; le glucose est une substance normalement absente de
l'urine, dont la présence signale un problème (souvent un diabète).

\corrige{de l'exercice 23}
Car l'hypertension abîme progressivement les vaisseaux sanguins des reins,
détériorant peu à peu leur capacité de filtration.

\corrige{de l'exercice 24}
Car elle évolue souvent silencieusement, sans symptôme visible, jusqu'à un stade déjà
avancé.

\corrige{de l'exercice 25}
Car lorsque les reins ne filtrent plus suffisamment le sang, les déchets et l'excès
d'eau et de sels s'accumulent dans l'organisme ; la dialyse les élimine
artificiellement à leur place.

\corrige{du sujet 1}
1) Reins, uretères, vessie, urètre.\\
2) Filtrer le sang puis réabsorber l'eau et les substances utiles.\\
3) $\dfrac{180-1{,}8}{180}\times100=99\,\%$.\\
4) Elles sont réabsorbées et renvoyées dans le sang.

\corrige{du sujet 2}
1) Le maintien constant de la composition du milieu intérieur de l'organisme.\\
2) Car le rein réabsorbe davantage d'eau pour compenser les pertes liées à la
transpiration.\\
3) Car les cellules ont besoin d'un environnement stable (eau, sels, déchets) pour
fonctionner normalement.\\
4) Les poumons, qui régulent notamment les taux de dioxygène et de dioxyde de carbone du
sang.

\corrige{du sujet 3}
1) Par exemple : le glucose, les protéines, le sang.\\
2) Un possible diabète.\\
3) Un problème de filtration rénale.\\
4) Car elle permet de détecter rapidement, simplement et sans douleur certains problèmes
de santé (diabète, infection, atteinte rénale).

\corrige{du sujet 4}
1) Une diminution de la capacité des reins à filtrer correctement le sang.\\
2) L'hypertension artérielle non contrôlée et le diabète mal équilibré.\\
3) Par exemple : fatigue, gonflements (œdèmes), diminution du volume d'urine.\\
4) Prévention : contrôler sa tension et son taux de glucose, éviter l'automédication
abusive. Traitement en cas de forme sévère : la dialyse (ou la greffe de rein).

\end{multicols}
\exercicesBanner
